%#########################
%			M A I N		D O C U M E N T
%#########################
%
% Here is where everything starts.
% Please read the README.md for general information on how this document works, folder-structer etc.
%
% Content of this file:
% Header:
% - SETTING UP DOCUMENT
% - IMPORTING DOCUMENT-INFO
% - SETTINGS
% 		- LOADING PACKAGES
% 		- SETTING UP PAGELAYOUT
% 		- SETTING UP GLOSSARY
% 		- SETTING UP BIBLATEX
% 		- PDF SETTINGS
% 		- SETUP GLOSSARY AND BIBLIOGRAPHY
% Body:
% - Content of the document
%
% via \usepackage{package name} settings are loaded! You only add your info in the document_info.sty file.
% If you want to include another chapter, appendix, etc. add it with \include{} to the body, next to the others.
%
%#########################
%			H E A D E R
%#########################
%
%--------------------------
%			S E T T I N G		UP		D O C U M E N T
%--------------------------
% this defines the general information of this document.
% look at this webside for all the options: https://texblog.org/2013/02/13/latex-documentclass-options-illustrated/
\documentclass[a4paper,		% Size of paper
               12pt,    	% Size of the text (fontsize)
               titlepage,       % Has titlepage
               twoside,	% Document is book styled
               DIV = 14,
               BCOR = 6mm,
               parskip = half,
               numbers = noenddot,
               headings	= small,
               headsepline = .4pt,
               cleardoublepage = empty]
               {scrbook} %scrbook is the type of document defined by KOMA
\newcommand{\mylanguage}{british}	%british or ngerman
%
%
%
%#########################
%			S E T T I N G S
%#########################
%
% You should not need to change the following files in the header.
%
%--------------------------
%			L O A D I N G		P A C K A G E S
%--------------------------
% this opens the packages.sty file which loads all the needed packages.
\usepackage{SETTINGS/packages_set}
\usepackage{SETTINGS/better_symbols}
\usepackage{SETTINGS/better_units}
%
%
%--------------------------
%			S E L E C T   L A N G U A G E   F O R   T I T E L P A G E
%--------------------------
% use "british" or "ngerman" 
\selectlanguage{\mylanguage}
% --------------------------
%		I M P O R T I N G		D O C U M E N T - I N F O
%--------------------------
% this loads all the general info on your document, like title etc.
% change this in the document_info.sty file.
% also add your own commands in there.
%
\usepackage{SETTINGS/document_info}
% --------------------------
%			S E T T I N G		P A G E L A Y O U T
%--------------------------
% this loads the pagelayout from layout.tex
\usepackage{SETTINGS/layout_set}
%
%--------------------------
%			S E T T I N G		U P		G L O S S A R I E S
%--------------------------
% this loads glossary.sty file and setups the glossary-style.
% Two glossaries will be created, one with latin and one with greek chars.
%
\usepackage{SETTINGS/glossary_set}
%
% TO ADD SOMETHING TO THE GLOSSARY USE Input_glossaries.tex FILE!
%
%
%--------------------------
% 			S E T T I N G		U P		B I B L A T E X
%--------------------------
%
% this loads all the settings for your bilbiography.
%
\usepackage{SETTINGS/biblatex_set}
%
%%--------------------------
% 			P D F   S E T T I N G S		H Y P E R R E F
%--------------------------
% this sets up the settings for the pdf, like filename.
% note: filename etc. are created from your settings in document_info.sty
%
% IMPORTANT: SET THE DRAFT SETTING TO TRUE FOR THE PRINT VERSION!
%
\usepackage{SETTINGS/pdf_set}
%
%
%--------------------------
% 			I N I T I A L I S I E R U N G			M A K E I N D E X
%--------------------------
%
\usepackage{makeidx}
\usepackage{float}
% Here I still do not fully understand why those are included here and cannot be
% included earlier...
% Kompatibilitaetsgruenden ganz zum Schluss einbinden

%
%--------------------------
% 			L O A D		G L O S S A R Y		/		 B I B L I O G R A P H Y
%--------------------------
\makeglossaries
\loadglsentries{BIB_AND_GLOSSARY/Input_glossaries} % *.tex filename of your glossaries.
% Standardmaessig werden im Glossar nur die Symbole angezeigt welche im Text mit \gls[}
% aufgerufen werden. Um alle definierten Eintraege anzuzeigen kann das \glsaddall Befehl
% verwendet werden:
% \glsaddall[types={notation,subscripts}]
\bibliography{BIB_AND_GLOSSARY/Biblio}	% *.bib filename of your literature.
%
%
%
% 
%
%
%
%
%
%
%
%##########################
%		B O D Y
%##########################
% The following contains the body of this document. With other words the content.
% This is where you actually start working. Add, change, remove chapters here.
% All your text is added with the \include{}.
% CHAPTERS_REST contains all the not so important files, like titlepage or summary.
% CHAPTERS_CONTENT is where you can put your chapters.
%
% Starts the body of the document
\begin{document}
% Defines footnotes
\deffootnote[1.5em]{1.5em}{1em}{\textsuperscript{\thefootnotemark}}
%
%--------------------------
% 			T I T L E P A G E
%--------------------------
% creates your titlepage
% Frontmatter activates roman pagenumbers. Pagestyle empty
\frontmatter
\pagestyle{empty}															% keine Kopf/Fusszeile, keine Seitenzahlen
%%% Titlepage
\begin{titlepage}
\begin{textblock*}{110mm}(88mm,12mm)
\includegraphics*[height=16mm]{./PICTURES/rwth_sla}
\end{textblock*}
\begin{textblock*}{\textwidth}(30mm,60mm)
\centering
{\large\bfseries \mythesistypeen} \\[3ex]
{\large\bfseries \mythesisnumber} \\[3ex]
\begin{huge}\bfseries\scshape \mytitleen\par\end{huge}
\end{textblock*}
\vspace*{85mm}
\centering\large
\phraseI \\[3ex]
\myinstitute\\
\myprof\\
\mydepart\\
\myuni\\[5ex]
\phraseII\\[1ex]
\textbf{\myname ~ \mysurname}\\
\vskip 5ex
\advisor: \\[1ex]
\myadvisorone \\
\myadvinstone\\[1ex]
%\begin{tabular}{p{.42\textwidth}p{.05\textwidth}p{.42\textwidth}}
%Dipl.-Ing. Max Musterbetreuer						& $\:$	&	Dr.-Ing. Emil Externbetreuer\\
%{\small Institut für Strukturmechanik und Leichtbau}&		& 	{\small Abteilung für Schnickschnack}\\[-.5ex]
%{\small RWTH Aachen	University}						& 		&	{\small Musterunternehmen, Musterstadt}
%\end{tabular}
\begin{textblock*}{\textwidth}(30mm,267mm)
\centering\large \myplace, \mydeadlineMonth ~ \mydeadlineYYYY
\end{textblock*}

\end{titlepage}
%% Titelseite Ende			\cleardoublepage
% Alternatively there seems to be an PDF file (by  Matsen) containing a spacial titlepage.
% You can add it with this command:
%\includepdf{titelseite.pdf}	\cleardoublepage
%
%--------------------------
% 			S U B J E C T			D E C L A R A T I O N   O N   T H E   U S E   O F   A I - B A S E D   T O O L S 
%--------------------------
%
% It has to be in german.
\pagestyle{empty}	% keine Kopf/Fusszeile, keine Seitenzahlen
\selectlanguage{ngerman}												% Erzwungen: Deutsch fuer Eigenstaendigkeitserklaerung
%% Eigenständigkeit. Neu als Erklärung über den Einsatz von generativen KI-Modellem
% Das Zentrale Prüfungsamt fordert folgende Datei: https://www.rwth-aachen.de/global/show_document.asp?id=aaaaaaaaaaoasiu
\chapter*{Erklärung über den Einsatz von generativen KI-Modellen in der Arbeit} 

Hiermit versichere ich, dass ich KI-gestützte Werkzeuge zu Text- oder Bildgenerierung in der vorliegenden Arbeit ausschließlich in den in der folgenden Tabelle gekennzeichneten Bereichen zu den jeweils angegebenen Zwecken verwendet habe. Ausgenommen von dieser Erklärung ist der Einsatz von KI-gestützten Werkzeugen für Zwecke, die sich nicht auf den wissenschaftlichen Inhalt der Arbeit auswirken (bspw. Grammatik-, Stil-, Rechtschreibprüfung und  Übersetzungsprogramme).

\vspace*{1cm}
\begin{minipage}{\textwidth}
\centering
\begin{tabularx}{\textwidth-4em}{Xl}
\toprule	
Abschnitt der Arbeit & Einsatzzweck \\ 
\midrule	
\textit{Seitenzahlen/ Unterkapitel / Abbildungen}& \textit{Art der Bearbeitung}\\%Hier werden Seitenzahlen/ Unterkapitel / Abbildungen genannt, die mit KI-gestützten Tools erstellt/bearbeitet wurden und die Art der Bearbeitung genannt. 		
\bottomrule	
\end{tabularx}
\end{minipage}

\vspace*{1cm}
Mir ist bekannt, dass ich dazu verpflichtet bin, Ergebnisse von generativen KI-Modellen immer gründlich zu prüfen. Ich übernehme die Verantwortung für alle Ergebnisse, die ich in meine Arbeit einschließe. Weiterhin ist mir bewusst, dass die ungeprüfte Übernahme von Halluzinationen einer generativen KI in meine Arbeit als Täuschungsversuch interpretiert werden muss.

%I hereby assure that I composed the thesis at hand self-reliantly.
\selectlanguage{ngerman}
\vspace*{1cm}\begin{tabular}{p{7cm}p{5cm}}
Aachen, den \today & $\rule{5cm}{0.01cm}$ \\
& (Unterschrift)
\end{tabular}
	\cleardoublepage							% Eigenstaendigkeitserklaerung
%
%--------------------------
% 			PREFACE + ABSTRACT
%--------------------------
\pagestyle{scrheadings} % Koma-script Kopf-/Fusszeile , Definition s.o.
%\pagenumbering{roman}														% Seitennummerierung in kleinen roemischen Zahlen
\begin{otherlanguage}{ngerman}
	\chapter*{Kurzfassung} 
\markboth{Kurzfassung}{Kurzfassung}
%% ==============================
Die Kurzfassung (Englisch abstract) hat das Ziel, den Leser über die wesentlichen Aspekte der Arbeit zu informieren. Hierzu gehören die Problemstellung, die Zielsetzung, der Lösungsweg sowie die wichtigsten Erkenntnisse der Arbeit. Dies alles soll in kurzer prägnanter Form auf maximal einer DIN A4-Seite zusammengefasst werden und keine Formeln oder Literaturhinweise enthalten. Dem Leser soll schnell erkennbar sein, was ihn beim Lesen der Arbeit erwartet. Dabei soll er feststellen können, was diese spezielle Arbeit ausmacht und inwiefern sie sich von anderen Arbeiten auf dem gleichen Gebiet unterscheidet.

Zusätzlich können ein bis zwei Bilder hinzugfügt werden, sofern sie zum wesentlichen Verständnis der Arbeit beitragen (z.B. Flussdiagramm). 

\begin{keywords}
Ca. fünf Schlagwörter zur treffenden Kennzeichnung des Dokumenteninhalts in alphabetischer Reihenfolge. Die Schlagwörter dienen als erste schnelle Entscheidungshilfe ob die Arbeit für den Leser inhaltlich relevant ist.\\
Z.B. Finite Elemente Methode, Faserverbundwerkstoffe, Stabilitätsverhalten, Zugversuch, Bruchmechanik.
\end{keywords} 	\cleardoublepage
\end{otherlanguage}
\begin{otherlanguage}{british}
	\chapter*{Abstract} 
\markboth{Abstract}{Abstract}
%% ==============================
Englische Version der Kurzfassung. Es gelten die gleichen Regeln.

Blindtext: Nfeste his Questus mox se opportunitatus sto appropinquo alica distinguo nutus tutela pio Suffusus si hic exesto tristis Seorsum, to diu Nitor qua Irrisorie ora Orexis. Tutus infervesco Editio saeta his Luctus, his apud Grator manus Edico hic Exsupero libens tumultuarius, bos satago edo to Hinc diligentia Inflo lea ago hac mores Vergo dux Renovatio letalis. No Declino vir excito utercumque Percut.

\begin{keywords}
\LaTeX, Keyword sample, Master thesis.
\end{keywords} \cleardoublepage
\end{otherlanguage}
% Falls ein Vorwort verwendet wird, entkommentiert die folgende Zeile
% \addchap*{Vorwort}
Sofern erforderlich, kann an dieser Stelle beispielsweise eine Danksagung erfolgen. Dies gilt für externe Arbeiten, ist bei internen Arbeiten jedoch nicht zulässig. Wenn kein Vorwort erforderlich ist, entfällt dieses Kapitel.
 \cleardoublepage
%
%--------------------------
%			S E L E C T   L A N G U A G E   F O R   L I S T I N G S
%--------------------------
% use "british" or "ngerman" 
\selectlanguage{\mylanguage}
% --------------------------
% 			LISTINGS
%    --------------------------
\tableofcontents \cleardoublepage
%	\listoffigures 				\cleardoublepage  % Wollte der Chef streichen
%	\listoftables				\cleardoublepage  % Wollte der Chef streichen
\printglossary[	type = \acronymtype,
		style = acrolongcol,
		title = Abkürzungsverzeichnis,
		nonumberlist] \cleardoublepage
\printglossary[	type = notation,
		style = symlongcolgrskip,
		nonumberlist] \clearpage
\printglossary[	type = subscripts,
        	style = symlongcol,
		nonumberlist] \cleardoublepage
%
%
%--------------------------
% 			CONTENT
%--------------------------
\mainmatter
\pagestyle{scrheadings}	% Koma-script Kopf-/Fusszeile , Definition s.o.
% Here you add chapters. Those are in the directory CHAPTERS_CONTENT.
% Numbers are only for sorting purpose in the directory.
\chapter{Einleitung}
\label{chap:Intro}
Die Einleitung liefert eine kurze Einordnung in den Projekthintergrund, um dem Leser den Einstieg in die Problemstellung zu ermöglichen. Sie kann je nach Art der Arbeit (Literaturrecherche, Konstruktion, praktische Arbeit, etc.) mehr oder weniger ausführlich sein, als Richtwert gilt ca. ein bis zwei Seiten.

Zu den Inhalten der Einleitung zählen:
%
\minisec{Zielsetzung}
An dieser Stelle erfolgt die Beschreibung der Aufgabenstellung, Motivation und Zielsetzung. Hier wird zunächst das zugrunde liegende Problem bzw. die zugrunde liegende Beobachtung erläutert. Anschließend wird auf die wichtigste Fragestellung bzw. die eigentliche Forschungsfrage eingegangen, woraus sich letztendlich das Ziel und die praktischen Bedeutung der Arbeit ableitet.
%
\minisec{Stand der Technik}
Hier wird \textbf{kurz} darauf eingegangen, wo die Aufgabenstellung im aktuellen Stand der Technik einzuordnen ist. Dabei werden sowohl bereits geleistete Vorarbeiten als auch vergleichbare Tätigkeiten referenziert. Zusätzlich wird hier deutlich gemacht, was den Neuheitsgrad dieser Arbeit ausmacht.
In diesem Kapitel wird zudem kurz umrissen, welche Methoden bisher bekannt sind, üblicherweise genutzt werden und inwieweit diese für die gegebene Aufgabenstellung von Relevanz sind. Dabei kann bereits eine Bewertung hinsichtlich der Zielsetzung erfolgen.
%
\minisec{Übersicht}
Zusätzlich soll die Einleitung eine kurze Übersicht über Aufbau und Gliederung der Arbeit geben in der kurz beschrieben wird, was in den einzelnen Kapiteln bearbeitet wird. Dies dient der Veranschaulichung des Roten Fadens dieser Arbeit.

Die einzelnen Unterpunkte können in einem Fließtext zusammengefasst werden.
		\cleardoublepage
\chapter{Stand der Technik}
\label{chap:SdT}

In diesem Kapitel soll der Stand der Technik dargestellt werden. Dieser soll sich insbesondere auf die Forschungshypothese bzw. Zielsetzung der Arbeit beziehen. Es soll aufgezeigt werden, auf welchem Stand sich die die Wissenschaft zum eigenen Forschungsthema befindet und wo das Potenzial liegt, die die eigene Arbeit rechtfertigt.		\cleardoublepage
\chapter{Methodik / Entwicklung der Methodik}
\label{chap:Principles}
Der Hauptteil der Arbeit gliedert sich in zwei Teile. Im ersten Teil (Methodik bzw. Entwicklung der Methodik) wird die grundlegende Methodik der Arbeit erarbeitet. Dabei werden grundlegende Fragen erläutert (Wie gehe ich vor? Was ist bekannt? Welche Annahmen werden getroffen? etc.), wobei die übergeordnete Forschungshypothese (bei einer MA) bzw. Zielsetzung der Arbeit (bei einer BA, PA, MT) berücksichtigt wird. Eine Erläuterung der Grundlagen ist nur da erlaubt, wo sie die täglichen Arbeiten eines Ingenieury der Strukturmechanik übersteigt. Der Umfang des Fachwissen kann sich in der Auswahl der Methodik, Interpretation der Ergebnisse und der kritischen Reflexion der gewählten Vorgehensweise widerfinden. Dieses Kapitel stellt die Grundlage der gesamten Arbeit dar.\\
Hier gelten zwei Regeln:
\begin{itemize}
	\item Von Experten - Für Experten!
	\item So wenig wie möglich und so viel wie nötig!
\end{itemize}

Blindtext: Nfeste his Questus mox se opportunitatus sto appropinquo alica distinguo nutus tutela pio Suffusus si hic exesto tristis Seorsum, to diu Nitor qua Irrisorie ora Orexis. Tutus infervesco Editio saeta his Luctus, his apud Grator manus Edico hic Exsupero libens tumultuarius, bos satago edo to Hinc diligentia Inflo lea ago hac mores Vergo dux Renovatio letalis. No Declino vir excito utercumque Percut.
\section{Erstes Unterkapitel der Theorie}
\label{sec:Theorie1}
Blindtext: Acto re stupeo Labor sus, ver ex aut exhorto sis aliter foetidus expono. Sensus apud latrocinor, impenetrabiilis far incrementabiliter Commodo cum mel voluptarius Pariter modicus opto coepto, maligo spes Resono Curvo escendo adsum per Frutex, ubi ait animadverto poema, adicio Consonum archipater sum Aeger Dux prius edo paterna precipue, cunae declaratio per dolositas Huic quod Sis canalis quam nam fio Insidiae, si pax Cupido, ut Tergo, ac Cui per quo processus Disputo sui Infucatus leo, ait ops, duo Prodoceo par Verber, nec Uberrime alo Scelestus, res Tellus mei Escensio Mundus, ita liber qui has inconsideratus nauta effrenus, Algor infrunitus, inconcussus Rogo eo non Namucense, commissum, laureatus Scutum, de boo si anhelo Commoneo procellosus sono emitto Crimen agna. Si subo Accubo castimonia hic ibi qua lux sto eu Pulcher Sem. Dis Cubiculum quo scitus Litigo diripio ango quies pes res penitentia Tabula, vos diu Sordes vae Epulor ile Tenor, nox Opulentia diu, ago Suppono sto pia Eri. \gls{FCS}			\cleardoublepage
\chapter{Umsetzung der Methodik}
\label{chap:Durchfuehrung}
Im zweiten Teil des Hauptteils folgt die Umsetzung der zuvor erarbeiteten Methodik. Aufbauend auf den im vorigen Kapitel erarbeiteten Grundlagen  wird nun der Lösungsweg aufgezeigt. Dabei wird das prinzipielle Vorgehen zum Erreichen der Zielsetzung unter Einbindung der verwendeten Hilfsmittel (Maschinen, Geräte, Programme etc. inklusive der verwendeten Einstellungen) aufgezeigt. Dabei müssen sämtliche verwendeten Daten ersichtlich sein, sodass es jederzeit möglich ist, die ermittelten Ergebnisse zu reproduzieren. 		\cleardoublepage
\chapter[Analyse und Auswertung der Ergebnisse]{Analyse und Auswertung der Ergebnisse - diese Kapitelüberschrift ist für das Inhaltsverzeichnis zu lang}
\label{chap:Implementation}
%Der Alternativtext in den eckigen Klammern ist die gekürzte Version für das Inhaltsverzeichnis

An dieser Stelle erfolgt eine Vorstellung der Ergebnisse. Hierbei sollen Randbedingungen berücksichtigt sowie eine kritische Analyse der Ergebnisse mit möglichen Unsicherheiten und Fehlern durchgeführt werden. 


Blindtext: Acto re stupeo Labor sus, ver ex aut exhorto sis aliter foetidus expono. Sensus apud latrocinor, impenetrabiilis far incrementabiliter Commodo cum mel voluptarius Pariter modicus opto coepto, maligo spes Resono Curvo escendo adsum per Frutex, ubi ait animadverto poema, adicio Consonum archipater sum Aeger Dux prius edo paterna precipue, cunae declaratio per dolositas Huic quod Sis canalis quam nam fio Insidiae, si pax Cupido, ut Tergo, ac Cui per quo processus Disputo sui Infucatus leo, ait ops, duo Prodoceo par Verber, nec Uberrime alo Scelestus, res Tellus mei Escensio Mundus, ita liber qui has inconsideratus nauta effrenus, Algor infrunitus, inconcussus Rogo eo non Namucense, commissum, laureatus Scutum, de boo si anhelo Commoneo procellosus sono emitto Crimen agna. Si subo Accubo castimonia hic ibi qua lux sto eu Pulcher Sem. Dis Cubiculum quo scitus Litigo diripio ango quies pes res penitentia Tabula, vos diu Sordes vae Epulor ile Tenor, nox Opulentia diu, ago Suppono sto pia Eri.

Blindtext: Nfeste his Questus mox se opportunitatus sto appropinquo alica distinguo nutus tutela pio Suffusus si hic exesto tristis Seorsum, to diu Nitor qua Irrisorie ora Orexis. Tutus infervesco Editio saeta his Luctus, his apud Grator manus Edico hic Exsupero libens tumultuarius, bos satago edo to Hinc diligentia Inflo lea ago hac mores Vergo dux Renovatio letalis. No Declino vir excito utercumque Percut.

Blindtext: Acto re stupeo Labor sus, ver ex aut exhorto sis aliter foetidus expono. Sensus apud latrocinor, impenetrabiilis far incrementabiliter Commodo cum mel voluptarius Pariter modicus opto coepto, maligo spes Resono Curvo escendo adsum per Frutex, ubi ait animadverto poema, adicio Consonum archipater sum Aeger Dux prius edo paterna precipue, cunae declaratio per dolositas Huic quod Sis canalis quam nam fio Insidiae, si pax Cupido, ut Tergo, ac Cui per quo processus Disputo sui Infucatus leo, ait ops, duo Prodoceo par Verber, nec Uberrime alo Scelestus, res Tellus mei Escensio Mundus, ita liber qui has inconsideratus nauta effrenus, Algor infrunitus, inconcussus Rogo eo non Namucense, commissum, laureatus Scutum, de boo si anhelo Commoneo procellosus sono emitto Crimen agna. Si subo Accubo castimonia hic ibi qua lux sto eu Pulcher Sem. Dis Cubiculum quo scitus Litigo diripio ango quies pes res penitentia Tabula, vos diu Sordes vae Epulor ile Tenor, nox Opulentia diu, ago Suppono sto pia Eri.

Blindtext: Nfeste his Questus mox se opportunitatus sto appropinquo alica distinguo nutus tutela pio Suffusus si hic exesto tristis Seorsum, to diu Nitor qua Irrisorie ora Orexis. Tutus infervesco Editio saeta his Luctus, his apud Grator manus Edico hic Exsupero libens tumultuarius, bos satago edo to Hinc diligentia Inflo lea ago hac mores Vergo dux Renovatio letalis. No Declino vir excito utercumque Percut.

Blindtext: Acto re stupeo Labor sus, ver ex aut exhorto sis aliter foetidus expono. Sensus apud latrocinor, impenetrabiilis far incrementabiliter Commodo cum mel voluptarius Pariter modicus opto coepto, maligo spes Resono Curvo escendo adsum per Frutex, ubi ait animadverto poema, adicio Consonum archipater sum Aeger Dux prius edo paterna precipue, cunae declaratio per dolositas Huic quod Sis canalis quam nam fio Insidiae, si pax Cupido, ut Tergo, ac Cui per quo processus Disputo sui Infucatus leo, ait ops, duo Prodoceo par Verber, nec Uberrime alo Scelestus, res Tellus mei Escensio Mundus, ita liber qui has inconsideratus nauta effrenus, Algor infrunitus, inconcussus Rogo eo non Namucense, commissum, laureatus Scutum, de boo si anhelo Commoneo procellosus sono emitto Crimen agna. Si subo Accubo castimonia hic ibi qua lux sto eu Pulcher Sem. Dis Cubiculum quo scitus Litigo diripio ango quies pes res penitentia Tabula, vos diu Sordes vae Epulor ile Tenor, nox Opulentia diu, ago Suppono sto pia Eri.

Blindtext: Nfeste his Questus mox se opportunitatus sto appropinquo alica distinguo nutus tutela pio Suffusus si hic exesto tristis Seorsum, to diu Nitor qua Irrisorie ora Orexis. Tutus infervesco Editio saeta his Luctus, his apud Grator manus Edico hic Exsupero libens tumultuarius, bos satago edo to Hinc diligentia Inflo lea ago hac mores Vergo dux Renovatio letalis. No Declino vir excito utercumque Percut.

Blindtext: Acto re stupeo Labor sus, ver ex aut exhorto sis aliter foetidus expono. Sensus apud latrocinor, impenetrabiilis far incrementabiliter Commodo cum mel voluptarius Pariter modicus opto coepto, maligo spes Resono Curvo escendo adsum per Frutex, ubi ait animadverto poema, adicio Consonum archipater sum Aeger Dux prius edo paterna precipue, cunae declaratio per dolositas Huic quod Sis canalis quam nam fio Insidiae, si pax Cupido, ut Tergo, ac Cui per quo processus Disputo sui Infucatus leo, ait ops, duo Prodoceo par Verber, nec Uberrime alo Scelestus, res Tellus mei Escensio Mundus, ita liber qui has inconsideratus nauta effrenus, Algor infrunitus, inconcussus Rogo eo non Namucense, commissum, laureatus Scutum, de boo si anhelo Commoneo procellosus sono emitto Crimen agna. Si subo Accubo castimonia hic ibi qua lux sto eu Pulcher Sem. Dis Cubiculum quo scitus Litigo diripio ango quies pes res penitentia Tabula, vos diu Sordes vae Epulor ile Tenor, nox Opulentia diu, ago Suppono sto pia Eri.

Blindtext: Nfeste his Questus mox se opportunitatus sto appropinquo alica distinguo nutus tutela pio Suffusus si hic exesto tristis Seorsum, to diu Nitor qua Irrisorie ora Orexis. Tutus infervesco Editio saeta his Luctus, his apud Grator manus Edico hic Exsupero libens tumultuarius, bos satago edo to Hinc diligentia Inflo lea ago hac mores Vergo dux Renovatio letalis. No Declino vir excito utercumque Percut.

Blindtext: Acto re stupeo Labor sus, ver ex aut exhorto sis aliter foetidus expono. Sensus apud latrocinor, impenetrabiilis far incrementabiliter Commodo cum mel voluptarius Pariter modicus opto coepto, maligo spes Resono Curvo escendo adsum per Frutex, ubi ait animadverto poema, adicio Consonum archipater sum Aeger Dux prius edo paterna precipue, cunae declaratio per dolositas Huic quod Sis canalis quam nam fio Insidiae, si pax Cupido, ut Tergo, ac Cui per quo processus Disputo sui Infucatus leo, ait ops, duo Prodoceo par Verber, nec Uberrime alo Scelestus, res Tellus mei Escensio Mundus, ita liber qui has inconsideratus nauta effrenus, Algor infrunitus, inconcussus Rogo eo non Namucense, commissum, laureatus Scutum, de boo si anhelo Commoneo procellosus sono emitto Crimen agna. Si subo Accubo castimonia hic ibi qua lux sto eu Pulcher Sem. Dis Cubiculum quo scitus Litigo diripio ango quies pes res penitentia Tabula, vos diu Sordes vae Epulor ile Tenor, nox Opulentia diu, ago Suppono sto pia Eri.

Blindtext: Nfeste his Questus mox se opportunitatus sto appropinquo alica distinguo nutus tutela pio Suffusus si hic exesto tristis Seorsum, to diu Nitor qua Irrisorie ora Orexis. Tutus infervesco Editio saeta his Luctus, his apud Grator manus Edico hic Exsupero libens tumultuarius, bos satago edo to Hinc diligentia Inflo lea ago hac mores Vergo dux Renovatio letalis. No Declino vir excito utercumque Percut.

Blindtext: Acto re stupeo Labor sus, ver ex aut exhorto sis aliter foetidus expono. Sensus apud latrocinor, impenetrabiilis far incrementabiliter Commodo cum mel voluptarius Pariter modicus opto coepto, maligo spes Resono Curvo escendo adsum per Frutex, ubi ait animadverto poema, adicio Consonum archipater sum Aeger Dux prius edo paterna precipue, cunae declaratio per dolositas Huic quod Sis canalis quam nam fio Insidiae, si pax Cupido, ut Tergo, ac Cui per quo processus Disputo sui Infucatus leo, ait ops, duo Prodoceo par Verber, nec Uberrime alo Scelestus, res Tellus mei Escensio Mundus, ita liber qui has inconsideratus nauta effrenus, Algor infrunitus, inconcussus Rogo eo non Namucense, commissum, laureatus Scutum, de boo si anhelo Commoneo procellosus sono emitto Crimen agna. Si subo Accubo castimonia hic ibi qua lux sto eu Pulcher Sem. Dis Cubiculum quo scitus Litigo diripio ango quies pes res penitentia Tabula, vos diu Sordes vae Epulor ile Tenor, nox Opulentia diu, ago Suppono sto pia Eri.

Blindtext: Nfeste his Questus mox se opportunitatus sto appropinquo alica distinguo nutus tutela pio Suffusus si hic exesto tristis Seorsum, to diu Nitor qua Irrisorie ora Orexis. Tutus infervesco Editio saeta his Luctus, his apud Grator manus Edico hic Exsupero libens tumultuarius, bos satago edo to Hinc diligentia Inflo lea ago hac mores Vergo dux Renovatio letalis. No Declino vir excito utercumque Percut.

Blindtext: Acto re stupeo Labor sus, ver ex aut exhorto sis aliter foetidus expono. Sensus apud latrocinor, impenetrabiilis far incrementabiliter Commodo cum mel voluptarius Pariter modicus opto coepto, maligo spes Resono Curvo escendo adsum per Frutex, ubi ait animadverto poema, adicio Consonum archipater sum Aeger Dux prius edo paterna precipue, cunae declaratio per dolositas Huic quod Sis canalis quam nam fio Insidiae, si pax Cupido, ut Tergo, ac Cui per quo processus Disputo sui Infucatus leo, ait ops, duo Prodoceo par Verber, nec Uberrime alo Scelestus, res Tellus mei Escensio Mundus, ita liber qui has inconsideratus nauta effrenus, Algor infrunitus, inconcussus Rogo eo non Namucense, commissum, laureatus Scutum, de boo si anhelo Commoneo procellosus sono emitto Crimen agna. Si subo Accubo castimonia hic ibi qua lux sto eu Pulcher Sem. Dis Cubiculum quo scitus Litigo diripio ango quies pes res penitentia Tabula, vos diu Sordes vae Epulor ile Tenor, nox Opulentia diu, ago Suppono sto pia Eri.

Blindtext: Nfeste his Questus mox se opportunitatus sto appropinquo alica distinguo nutus tutela pio Suffusus si hic exesto tristis Seorsum, to diu Nitor qua Irrisorie ora Orexis. Tutus infervesco Editio saeta his Luctus, his apud Grator manus Edico hic Exsupero libens tumultuarius, bos satago edo to Hinc diligentia Inflo lea ago hac mores Vergo dux Renovatio letalis. No Declino vir excito utercumque Percut.


	\cleardoublepage
\chapter{Schlussfolgerung und Ausblick}
\label{Conclusion}
Zum Schluss wird die Arbeit rekapituliert, insbesondere im Hinblick auf die Forschungshypothese (MA) bzw. Zielsetzung der Arbeit (BA).
Mögliche Fragestellungen zur Orientierung:
\begin{itemize}
	\item Was wurde in dieser Arbeit gemacht?
	\item Was sind die wesentlichen Ergebnisse?
	\item Konnte die Forschungshypothese unterstützt oder falsifiziert werden? (Analog Zielsetzung der Arbeit: Konnte sie erreicht werden?) Wenn nein, warum nicht?
	\item Wo ist weiterer Forschungsbedarf?
	\item Welche interessanten Forschungsbereiche ergeben sich aus der eigenen Arbeit?
\end{itemize} 	\cleardoublepage
%
%--------------------------
% 			BIBLIOGRAPHY
%--------------------------
\addcontentsline{toc}{chapter}{Bibliography}% add literature to  table of contents.
\nocite{*} % Alle Literatureintraege ins Vz. aufnehmen
% BITTE AUSKOMMENTIEREN FUER DIE ARBEIT
% NUR ZU ANSCHAUUNGSZWECKEN IN DEM BEISPIEL
\printbibliography \cleardoublepage
%
%--------------------------
% 			APPENDIX
%--------------------------
\appendix
\pagestyle{scrheadings}	% Koma-script Kopf-/Fusszeile , Definition s.o.
\chapter{Anhang - Tipps zum Verfassen studentischer Arbeiten}
\minisec{Allgemeine Regeln:}
\begin{compactitem}
\item[\emph{Neutrale Formulierung:}] Die Arbeit wird in der 3. Person Singular verfasst.
\item[\emph{Tempus:}] Die Arbeit wird durchgängig im Präsens geschrieben.
\item[\emph{Aktive Formulierung:}] Formulierungen im Passiv vermeiden.
\item[\emph{Formulierung im Indikativ:}] Die Verwendung des Konjunktivs vermeiden.
\item[\emph{Abkürzungen:}] Alle Abkürzungen, die nicht im Duden stehen, müssen eingeführt werden.
\end{compactitem}
\minisec{Generell gilt:}
\begin{itemize}
\item[\emph{So wenig wie möglich, so viel wie nötig!}] Herleitungen und Erklärungen auf das Wesentliche beschränken und nicht zu weit abschweifen; Vermeidung der Aufzählung unnötiger Einzelheiten, von allgemein Bekanntem und von Informationen, die nicht zum Thema gehören.
\item[\emph{So einfach wie möglich:}] Möglichst kurz und prägnant formulieren. 
\item[\emph{Kurze Sätze:}] Vermeidung von komplizierten Satzreihen und Schachtelsätzen sowie von langen zusammengesetzten Wörtern.
\item[\emph{Fachvokabular Sätze:}] An den richtigen Stellen einheitlich und durchgängig verwenden. Für das Verständnis notwendige, aber nicht geläufige Fach- und Fremdwörter sowie Abkürzungen erläutern.
\item[\emph{Redundanz:}] Unbekannte und besonders wichtige Informationen müssen dem Leser durch angemessene Wiederholung vors Auge geführt werden. 
\item[\emph{Übersichtlichkeit:}] Ergebnisse, Versuchspläne, Maschinendaten, etc. möglichst in Tabellenform abfassen und nicht im Fließtext aufzählen.
\item[\emph{Aussagekräftige Kapitelüberschriften}] (nicht zu lang!) erleichtern das Verständnis.
\item[\emph{Unterkapitel}] sind nur bis zur vierten Stufe zulässig; Eine neue Stufe ist nur dann einzuführen, wenn es mindestens zwei Kapitel in dieser Ebene gibt.
\item[\emph{Kapiteleinleitung und -abschluss:}] Ein bis zwei \emph{kurze Einleitungssätze} am Anfang von jedem Kapitel (z.B. In diesem Kapitel geht es um..., Im Folgenden wird beschrieben...) bzw. \emph{zusammenfassende Sätze} am Ende eines Kapitels geben dem Text eine Struktur und tragen so zum allgemeinen Verständnis bei.
\end{itemize}

Der Umfang der schriftlichen Arbeit richtet sich nach der jeweiligen Prüfungsordnung. Es gilt:
\begin{center}
\begin{tabular}{lr}
\toprule
{Projektarbeit} & k.A.\\
{Bachelorarbeit} & < 50 Seiten ohne Anhang\\
{Masterarbeit} & < 80 Seiten ohne Anhang\\
\bottomrule
\end{tabular}
\end{center}
%
\minisec{Kontrollfragen:}
\begin{itemize}
\item Ist der \emph{Rote Faden} durch die ganze Arbeit erkennbar?
\item Folgt der Aufbau der Arbeit einer \emph{logischen Struktur} (Reihenfolge, Gedankensprünge, Sinnabschnitte,etc.)?
\item Gibt es eine durchgängige \emph{Zeitform} (Präsens), oder gibt es Zeitsprünge (z.B. Präteritum im Durchführungsteil)?
\item Stimmen die \emph{inhaltlichen Bezüge}? (Passt die Verbform zum Subjekt?)
\item Sind die Inhalte \emph{übersichtlich} gestaltet?
\end{itemize}


 \cleardoublepage
\chapter{Anhang - Formale Fragen}
Die Formatierung der Vorlage ist beizubehalten. Die Farben in Diagrammen und Zeichnungen müssen so gewählt werden, dass durch das Ausdrucken der Arbeit in Schwarz-Weiß keine Informationen verloren gehen (Graustufen). Die Verwendung der Farbe gelb ist unzulässig. Weiterhin muss auf Personen mit einer Rotgrünblindheit bei der Farbwahl Rücksicht genommen werden. Die Zahlen von eins bis einschließlich zwölf werden ausgeschrieben. Höhere Zahlen werden \emph{nicht} ausgeschrieben. Diagramme, Tabelle, etc. sind so zu formatieren, dass sie inklusive ihrer Unterschriften bei einer Verkleinerung auf DIN A5 lesbar bleiben.
%
\section{Umgang mit Latex - Beispiele}
%
\minisec{Formeln}
Die Theorie beinhaltet Formeln, Herleitungen, empirische Daten, etc. Für die Darstellung von Formeln in der Formelumgebung bzw. im Fließtext gelten einige Regeln:

Das Multiplikationszeichen ist der Punkt und nicht der Stern. Verwendet wird er bei Skalarprodukten, oder wenn sonst Irrtümer entstehen können. Ansonsten macht ein Multipliaktionszeichen Formeln unübersichtlich lang. Um die Eindeutigkeit und Übersichtlichkeit in Formeln zu gewährleisten bietet es sich häufig an zusätzliche Abstandshalter und verschieden große Klammern zu verwenden.
\begin{equation}
\begin{split}
%\dot{\gls{x}}&~=~\gls{A}\gls{x}(t)~+~%
%							  \gls{B}\gls{u}(t)\\
%\gls{y}&~=~\gls{C}\gls{x}(t)~+~\gls{D}\gls{u}(t)\\
\dot{\underline{x}}&~=~\underline{A}\underline{x}(t)~+~%
							  \underline{B}\underline{u}(t)\\
\underline{y}&~=~\underline{C}\underline{x}(t)~+~\underline{D}\underline{u}(t)\:.
\end{split}
\label{eq:Gleichung1}
\end{equation} \glsadd{A}\glsadd{B}\glsadd{C}\glsadd{D}\glsadd{x}\glsadd{u}
\begin{equation}
\begin{pmatrix}
	\delta\dot{q_k}		\\
	\delta\dot{\alpha_k}	\\
	\delta\dot{V_k}		\\
	\delta\dot{\gamma_k} \end{pmatrix}%
=\underbrace{\begin{bmatrix}
	M_q			& M_{\alpha} 	& M_u 	& 0	\\
	1+Z_q 		& Z_{\alpha} 	& Z_u	& 0	\\
	X_q			& X_{\alpha}-g	& X_u 	& -g\\
	-Z_q 		& -Z_{\alpha} 	& -Z_u	& 0	\end{bmatrix}}_{\text{Zustandsmatrix}}%
\begin{pmatrix}
	\delta{q_k}			\\
	\delta{\alpha_k}	\\
	\delta{V_k}			\\
	\delta{\gamma_k} 	\end{pmatrix}%
+\underbrace{\begin{bmatrix}
	M_{F} 		& M_{\kappa}	& M_{\eta}\\
	Z_{F} 		& Z_{\kappa}	& Z_{\eta}\\
	X_{F} 		& X_{\kappa}	& X_{\eta}\\
	-Z_{F} 		& -Z_{\kappa}	& -Z_{\eta}\end{bmatrix}}_{\text{Eingangsmatrix}}%
\begin{pmatrix}
	\delta_F		\\
	\delta\kappa	\\
	\delta\eta 		\end{pmatrix}\:.
\label{eq:gleichung2}
\end{equation}
In der Formel-Umgebung wird Text so dargestellt, wie im Fließtext. Dazu verwendet man den Befehl $\backslash$text oder $\backslash$mathrm. Letzterer empfiehlt sich u.a. für Einheiten im Formeltext, da lediglich die Schriftart geändert wird, die Formel-Umgebung jedoch erhalten bleibt. Auch der Differentialoperator (Nicht zu verwechseln mit der partiellen Ableitung) bzw. Integrale werden nicht kursiv geschrieben:
\begin{equation}
E=\frac{\operatorname{d} \! \sigma}{\operatorname{d} \! \epsilon} \:.
\end{equation}
Partielle Ableitung:
\begin{equation}
\tau=G \;\left(\frac{\partial u}{\partial y} + \frac{\partial v}{\partial x}\right) \:.
\end{equation}

In abgesetzten Formeln werden Brüche mit horizontalem Bruchstrich (Befehl $\backslash$frac{}{}) verwendet  $\mathrm{\frac{N}{mm^2}}$. Im Fließtext werden Brüche mit einem Schrägstrich dargestellt. Dabei ist ausdrücklich die Formatierung \emph{ohne} Fraction-Umgebung, $\mathrm{N/mm^2}$, der Formatierung \emph{mit} einer Umgebung, wie dem Befehl $\backslash$sfrac{}{}, $\mathrm{\sfrac{N}{mm^2}}$, vorzuziehen.

Die Zeichen für die Dezimal- und Tausendertrennung richten sich nach der Sprache der Arbeit (vgl. Tabelle \ref{tab:Trennzeichen}). In Grafiken sollte man sich an die jeweiligen Konventionen halten, in deutschsprachigen Arbeiten ist es jedoch ebenso zulässig, die englische Konvention zu nutzen, sofern eine andere Darstellung nicht mit vertretbarem Aufwand erreichbar ist.
\begin{table}
\centering
\captionof{table}{Dezimal- und Tausendertrennung, aufgeführt nach Sprachraum.}
\begin{tabular}{lrr}
\toprule
							& Deutsch	& Englisch\\
\midrule
Tausendertrennzeichen		& .			& ,\\
Dezimaltrennzeichen 		& ,			& .\\
\bottomrule
\end{tabular}
\label{tab:Trennzeichen}
\end{table}
%
\minisec{Verwendung von Summenzeichen in mathematischen Formeln}
In mathematischen Berechnungen stehen sowohl das Summenzeichen ($\sum$) \texttt{(\textbackslash sum)} als auch das erweiterte Summenzeichen ($\Xsum$) \texttt{(\textbackslash Xsum)} zur Verfügung.

$$\dfrac{\pi^2}{6} = \sum_{n=1}^\infty \dfrac{1}{n^2}$$

Es wird jedoch empfohlen, das erweiterte Summenzeichen $\Xsum$ zu verwenden, da dies eine anschaulichere Darstellung ermöglicht. Daher ist die bevorzugte Notation:
$$\dfrac{\pi^2}{6} = \Xsum_{n=1}^\infty \dfrac{1}{n^2}$$

\minisec{Abbildungen}
Das Hinzufügen von Grafiken und Abbildungen dient der Veranschaulichung bzw. der Ergänzung des Fließtexts. Abbildungen werden mit der $\backslash$figure-Umgebung in den Text eingebunden. Der Befehl mit dem die jeweilige Abbildung aufgerufen wird lautet:\\
$\backslash$includegraphics[ scale=.71]\{pictures/Dateiname.eps\}. Dazu muss die verwendete Datei als .eps- bzw. als .png-Datei im Ordner \glqq pictures\grqq \! abgelegt werden. Mit Befehlen wie beispielsweise scale, angle, width, oder height in einer Latex-spezifischen Einheit wird das Bild bezüglich Skalierung, Drehwinkel, Breite oder Höhe angepasst. Bei Zeichnungen ist darauf zu achten, dass Zeichnungs- und Bemaßungslinien, sowie Schriftbreiten mit einem normgerechten Satz von Strichbreiten erstellt wurden. Bei Verwendung des Strichbreiten-Satzes [0,35mm; 0,7mm] ist in der Zeichnung eine Schriftgröße von 16pt in etwa normgerecht. Eine Skalierung dieser Zeichnung mit dem Faktor $\mathrm{=\:1/\sqrt{2}}$ ergibt einen Strichbreitensatz von [0,25mm; 0,5mm] und die entsprechende Schriftgröße von 12pt. Ein Beispiel ist in Abbildung \ref{fig:Bsp_Zeichung} dargestellt.
\begin{figure}[t!]
\centering
\includegraphics[scale=0.71]{PICTURES/Beispiel.eps}
\caption{Dieses Beispielbild wurde mit dem normgerechten Satz der Strichbreiten [0,35mm; 0,7mm] und einer Schriftgröße von 16pt erstellt und gemäß der gewählten Schriftgröße im Fließtext von 12pt mit dem erforderlichen Faktor von scale=.71 ($\mathrm{=\:1/\sqrt{2}}$) skaliert.}
\label{fig:Bsp_Zeichung}
\end{figure}
%
\minisec{Tabellen}
Tabelle \ref{tab:refCase} zeigt ein Beispiel für eine Standard-Tabelle mit der Umgebung \emph{$\backslash$tabular}. Hier wird auch die Nutzung einiger weiterer typischer Befehle in Tabellen deutlich. Sollte eine Tabelle von sich aus nicht über die gesamte Breite gehen, bietet es sich aus optischen Gründen häufig an, diese zu über eine definierte Textbreite zu strecken. Dies erfolgt mithilfe des Pakets TabularX. Ein Beispiel ist in Tabelle \ref{tab:refTabX} gezeigt. Sollte jedoch eine Tabelle so schmal ausfallen, dass die Streckung unvorteilhaft wird, so kann eine einfache Tabular-Umgebung, wie in Tabelle \ref{tab:refCase_simple} gewählt werden.

\begin{table}[htb]
	\centering
	\captionof{table}{Dies ist eine Tabellenüberschrift, die so lang ist, dass sie sich über zwei Zeilen erstreckt. Man erkennt, dass die nachfolgenden Zeilen eingerückt werden.}
	\label{tab:refCase}
	\begin{tabularx}{\textwidth-4em}{m{6mm}Xrrrrrrrr}
		\toprule
		&		& \multicolumn{8}{c}{Mach number}\\
		\cmidrule{3-10}
		& 		& 0.2	& 0.4	& 0.6	& 0.75	& 0.82	& 0.85	& 0.88& 0.96\\
		\midrule %\cmidrule{3-10}
		\multirow{7}{*}{\begin{sideways} dyn. pressure $\mathrm{[10^3 kg\:/m\:s^2]}$ \end{sideways}}%
		& 2.00 	& 2800	& 12000	& 0		& 0		& 0		& 0		& 0		& 0	\\
		& 2.30 	& 1700	& 11000	& 0		& 0		& 0		& 0		& 0		& 0	\\
		& 2.70 	& 400	& 10000	& 0		& 0		& 0		& 0		& 0		& 0	\\
		& 3.30 	&-1000	& 9300	& 14000	& 0		& 0		& 0		& 0		& 0	\\
		& 4.10 	& 0		& 7600	& 13000	& 1500	& 0		& 0		& 0		& 0	\\
		& 5.80 	& 0		& 5300	& 10000	& 1200	& 14000	& 15000	& 15000 & 0	\\
		& 8.00 	& 0		& 2800	& 8700	& 1100	& 12000	& 13000	& 13000 & 14000\\
		\midrule
		\multicolumn{10}{c}{altitude [$\mathrm{m}$]}\\
		\bottomrule
	\end{tabularx}
\end{table}

\begin{table}[htb]
	\centering
	\captionof{table}{Tabellenüberschrift}
	\begin{tabular}{lrr}
		\toprule
		Werkstoff	& Dichte				& E-Modul\\
		& $\mathrm{[g/cm^3]}$	& $\mathrm{[N/mm^2]}$\\
		\midrule
		Stahl		& 7,85					& 210000\\
		AlCuMg1 	& 2,80					& 71500\\
		CFK			& 1,50					& 150000\\
		\bottomrule
	\end{tabular}
	\label{tab:refCase_simple}
\end{table}

Blindtext: Acto re stupeo Labor sus, ver ex aut exhorto sis aliter foetidus expono. Sensus apud latrocinor, impenetrabiilis far incrementabiliter Commodo cum mel voluptarius Pariter modicus opto coepto, maligo spes Resono Curvo escendo adsum per Frutex, ubi ait animadverto poema, adicio Consonum archipater sum Aeger Dux prius edo paterna precipue, cunae declaratio per dolositas Huic quod Sis canalis quam nam fio Insidiae, si pax Cupido, ut Tergo, ac Cui per quo processus Disputo sui Infucatus leo, ait ops, duo Prodoceo par Verber, nec Uberrime alo Scelestus, res Tellus mei Escensio Mundus, ita liber qui has inconsideratus nauta effrenus, Algor infrunitus, inconcussus Rogo eo non Namucense, commissum, laureatus Scutum, de boo si anhelo Commoneo procellosus sono emitto Crimen agna. Si subo Accubo castimonia hic ibi qua lux sto eu Pulcher Sem. Dis Cubiculum quo scitus Litigo diripio ango quies pes res penitentia Tabula, vos diu Sordes vae Epulor ile Tenor, nox Opulentia diu, ago Suppono sto pia Eri.

\begin{table}[htb]
	\centering
	\captionof{table}{Dies ist eine Tabelle mit TabularX, deren Breite auf die Textbreite abzüglich je links und rechts 2em ausgestreckt wurde.}
	\begin{tabularx}{\textwidth-4em}{Xrrr}
		\toprule
		&$\omega_0~\mathrm{[rad/s]}$ &$\zeta~\mathrm{[-]}$& Polstellen\\
		\midrule
		Erste Biegemode - ungeregelt	& 10.0000	& 0.0700	& -0.6700 $\pm$ 10.0036i\\
		Erste Biegemode - geregelt		& 9.9000	& 0.0750	& -0.7500 $\pm$ 9.9652i\\[1ex]
		AS-Mode - ungeregelt			& 1.9000	& 0.5700	& -1.1000 $\pm$ 1.5600i\\
		AS-Mode - geregelt 				& 1.9000	& 0.5700	& -1.1000 $\pm$ 1.5600i \\[1ex]
		Phygoiden-Mode - ungeregelt		&			& 			& -0.2000 $\wedge$ 0.3000\\
		Phygoiden-Mode - geregelt		&			& 			& -0.2000 $\wedge$ 0.3000\\
		\bottomrule
	\end{tabularx}
	\label{tab:refTabX}
\end{table}
%
\minisec{Nützliche Befehle}
Mit dem Befehl $\backslash$emph werden wesentliche Begriffe \emph{hervorgehoben}. Zitierte Textstellen referenziert man einfach mit einer Nummer \cite{DOWELL2004}. Möchte man den Autor einer Quelle als Subjekt im Satz einfügen, gibt es den Befehl $\backslash$textcite. Ein Beispiel: \textcite{DOWELL2004} beschreiben wesentliche Aspekte.

Absätze werden, wie bekannt aus MS Word, mit einer Leerzeile erzeugt. Einen Zeilensprung erhält man mit dem Befehl $\backslash\backslash$. Zur Veranschaulichung folgt ein Beispiel. Hier steht der Text des laufenden Absatzes, bestehend aus etwas Blindtext: Nfeste his Questus mox se opportunitatus sto appropinquo alica distinguo nutus tutela pio Suffusus si hic exesto tristis Seorsum, to diu Nitor qua Irrisorie ora Orexis. Dieser Satz beendet den laufenden Absatz mit dem Befehl $\backslash\backslash$.\\
Dieser Satz folgt auf den einfachen Zeilensprung mit dem Befehl $\backslash\backslash$. Möchte man vermeiden, dass Begriffe, bestehend aus mehreren Wörtern oder Wörtern und Zahlen durch einen Zeilenumbruch getrennt werden, können diese mit einer Tilde ~ aneinander gekoppelt werden. Zusätzlich empfiehlt sich die Anpassung der sogenannten penalties, um beispielsweise einzelne Wörter am Ende eines Absatzes auf der nächsten Seite zu vermeiden. Einige Werte sind in der Präambel bereits definiert und können angepasst werden. Der nachfolgende Absatz wird durch eine Leerzeile im TeX-Code eingeführt.

Um bestimmte Punkte kurz aufzulisten empfiehlt sich die Umgebung \emph{compactitem}. Im Vergleich zu der Umgebung \emph{itemize} rückt dieser Befehl die einzelnen Auzählungspunkte näher zusammen:
\begin{compactitem}
\item Punkt 1 beinhaltet etwas Text,
\item Punkt 2 ist auch interessant und
\item Punkt 3 ist ebenfalls ein wesentlicher Aspekt.
\end{compactitem}

Vor der Einleitung wird das Symbol- und Abkürzungsverzeichnis erstellt. Um Abkürzungen oder Symbole im Symbolverzeichnis zu referenzieren verwendet man die Befehle $\backslash$gls, bzw. $\backslash$glsadd, falls im Verzeichnis nur ein Verweis auf diese Seite erfolgen soll. Der Befehl $\backslash$gls erzeugt unter Verwendung des hyperref packages (mit der Option: draft = false) einen markierten Link in das Symbolverzeichnis. Im Symbolverzeichnis beschriebene Begriffe erscheinen erst, wenn diese im Text mittels einer der Befehle referenziert werden. Dies kann vollständig mit dem Befehl $\backslash$glsaddall[types=\{nomenclature, subscripts\}] umgangen werden. \glsadd{ROMAN}\glsadd{GREEK} \glsadd{eta}\glsadd{Theta}\glsadd{gamma}\glsadd{kappa}\glsadd{alpha} \glsaddall[types={subscripts,notation}] %

Kompiliert wird der Text mit der Befehlsfolge pdflatex biber.exe pdflatex.exe. Um das Symbolverzeichnis mit dem Package \emph{glossaries} lauffähig zu machen, bitte in die Erläuterung der Präambel schauen.
 \cleardoublepage
%	\chapter{Anhang mit wesentlichem Code (Beispiel hier: Matlab)}
\label{chap:Anhang_Matlab}
\lstinputlisting{DATA/MATLAB_Plot_Vorlage_SLA_2015_11_11.m}
			\cleardoublepage
%
\end{document}
%
%
%
%% end of file, you have done it!
